% Copyright 2026 MarcosHCK
%
\documentclass{article}[12pt,a4paper]
\usepackage{amsfonts}
\usepackage{amsmath}
\usepackage{amsthm}
\usepackage{bookmark}
\usepackage{changepage}
\usepackage{hyperref}
\usepackage{cleveref} % must load after hyperref
\usepackage{float}
\usepackage{graphicx}
\usepackage{inputenc}[utf8]
\usepackage{indentfirst}
\usepackage{polyglossia}[spanish,english]
\usepackage{biblatex}[style=apa] % must load after polyglossia
\usepackage{csquotes} % must load after biblatex
\usepackage{setspace}
\usepackage{pgf}
\usepackage{xcolor}[dvipsnames]

\setmainlanguage{spanish}
\setotherlanguage{english}

\setstretch{1.5}

\def\mathdefault#1{#1}
\newcommand{\myeoch}{\pagebreak}

\definecolor{Aquamarine}{rgb}{0.498, 1, 0.83}
\definecolor{DarkAquamarine}{rgb}{0.18, 0.55, 0.34}
\definecolor{YellowOrange}{rgb}{1, 0.7, 0}
\newcommand\myshade{85}
\colorlet{mycitecolor}{YellowOrange}
\colorlet{mylinkcolor}{violet}
\colorlet{myurlcolor}{DarkAquamarine}

\hypersetup {
    colorlinks = true,
    citecolor  = mycitecolor!\myshade!black,
    linkcolor  = mylinkcolor!\myshade!black,
    urlcolor   = myurlcolor!\myshade!black,
  }

\begin{document}

\begin{titlepage}

  \centering
  {\huge\bfseries Reporte \#2\par}
	{\Large Autor: \itshape Marcos Antonio Pérez Lorenzo\par}
\end{titlepage}

\pagebreak

\section{Teoría}

Sea la función \(f(x)\) donde:

\begin{equation}
  f(x)=\begin{cases}
     3 & -\pi + 2k\pi < x < 0 + 2k\pi \\
    -7 & 0 + 2k\pi < x < \pi + 2k\pi
  \end{cases}
  \label{eq:fx}
\end{equation}

Donde \(k\in\mathbb{R}\). Ver figura \ref{fig:original}.

Se quiere analizar la convergencia de una aproximación mediante ciertas cantidades de elementos de la Serie de Fourier correspondiente a \(f(x)\). Para ello, primero es necesario calcular los coeficientes \(A_0\), \(A_n\) y \(B_n\):

\begin{proof}

  \begin{equation}
    f(x)=\frac{A_0}{2}+\sum_{n=1}^{\infty}\left[A_n\cos\left(\frac{2n\pi}{T}x\right)+B_n\sin\left(\frac{2n\pi}{T}x\right)\right]
    \label{eq:fourierseries}
  \end{equation}

  Sea \(P\) un intervalo definido como \(P=\left[-\frac{T}{2},\frac{T}{2}\right]\)

  \begin{equation}
    A_0=\frac{2}{T}\int_{P}f(x)\,dx
    \label{eq:a0coeff}
  \end{equation}
  \begin{equation}
    A_n=\frac{2}{T}\int_{P}f(x)\cos\left(\frac{2n\pi}{T}x\right)\,dx
    \label{eq:ancoeff}
  \end{equation}
  \begin{equation}
    B_n=\frac{2}{T}\int_{P}f(x)\sin\left(\frac{2n\pi}{T}x\right)\,dx
    \label{eq:bncoeff}
  \end{equation}

  Para la función \ref{eq:fx}, \(P=\left[-\pi,\pi\right]\), el coeficiente \(A_0\) queda:

  \begin{equation*}
    A_0=\frac{2}{T}\int_{P}f(x)\,dx=\frac{1}{\pi}\int_{-\pi}^{0}3\,dx-\frac{1}{T}\int_{0}^{\pi}7\,dx
  \end{equation*}
  \begin{equation}
    A_0=\frac{1}{\pi}\left[3x\right]_{-\pi}^{0}-\frac{1}{\pi}\left[7x\right]_{0}^{\pi}=-4
    \label{eq:a0coeffr}
  \end{equation}

  \(A_n\) se desarrolla de la forma:

  \begin{equation*}
    A_n=\frac{2}{T}\int_{P}f(x)\cos\left(\frac{2n\pi}{T}x\right)\,dx=\frac{3}{\pi}\int_{-\pi}^{0}\cos\left(nx\right)\,dx-\frac{7}{\pi}\int_{0}^{\pi}\cos\left(nx\right)\,dx
  \end{equation*}
  \begin{equation*}
    \int_{a}^{b}\cos\left(nx\right)\,dx=\left[\frac{\sin\left(nx\right)}{n}\right]_{a}^{b}
  \end{equation*}
  \begin{equation}
    A_n=\frac{3}{\pi}\left[\frac{\sin\left(nx\right)}{n}\right]_{-\pi}^{0}-\frac{7}{\pi}\left[\frac{\sin\left(nx\right)}{n}\right]_{0}^{\pi}=0
    \label{eq:ancoeffr}
  \end{equation}

  Finalmente \(B_n\) se expande como:

  \begin{equation*}
    B_n=\frac{2}{T}\int_{P}f(x)\sin\left(\frac{2n\pi}{T}x\right)\,dx=\frac{3}{\pi}\int_{-\pi}^{0}\sin\left(nx\right)\,dx-\frac{7}{\pi}\int_{0}^{\pi}\sin\left(nx\right)\,dx
  \end{equation*}
  \begin{equation*}
    \int_{a}^{b}\sin\left(nx\right)\,dx=\left[-\frac{\cos\left(nx\right)}{n}\right]_{a}^{b}
  \end{equation*}
  \begin{equation*}
    \int_{-\pi}^{0}\sin\left(nx\right)\,dx=\left[-\frac{\cos\left(nx\right)}{n}\right]_{-\pi}^{0}=\left[-\frac{1}{n}\right]-\left[-\frac{\cos\left(n\pi\right)}{n}\right]=\frac{\cos\left(n\pi\right)-1}{n}
  \end{equation*}
  \begin{equation*}
    \int_{0}^{\pi}\sin\left(nx\right)\,dx=\left[-\frac{\cos\left(nx\right)}{n}\right]_{0}^{\pi}=\left[-\frac{\cos\left(n\pi\right)}{n}\right]-\left[-\frac{1}{n}\right]=\frac{1-\cos\left(n\pi\right)}{n}
  \end{equation*}
  \begin{equation*}
    B_n=\frac{3}{\pi}\frac{\cos\left(n\pi\right)-1}{n}-\frac{7}{\pi}\frac{1-\cos\left(n\pi\right)}{n}=\frac{10}{\pi}\frac{\cos\left(n\pi\right)-1}{n}
  \end{equation*}

  Dado que \(cos\left(n\pi\right)=(-1)^n\), entonces obtenemos la expresión:

  \begin{equation}
    B_n=10\frac{(-1)^n-1}{n\pi}
    \label{eq:bncoeffr}
  \end{equation}

  Utilizando los resultados \ref{eq:a0coeffr}, \ref{eq:ancoeffr} y \ref{eq:bncoeffr} expandiéndolos en la expresión \ref{eq:fourierseries}, obtenemos:

  \begin{equation}
    f(x)=-2+\frac{10}{\pi}\sum_{k=0}^{\infty}\frac{\sin\left(\left(2k+1\right)x\right)}{2k+1}
  \end{equation}

\end{proof}

\section{Figuras}

\subsection{Función original}

\begin{figure}[H]
  \begin{adjustwidth}{-0.33\linewidth}{-0.33\linewidth}
    \resizebox{\linewidth}{!}{\input{fx.pgf}}
  \end{adjustwidth}
  \caption{f(x) original}
  \label{fig:original}
\end{figure}

\subsection{Función aproximada}

\begin{figure}[H]
  \begin{adjustwidth}{-0.33\linewidth}{-0.33\linewidth}
    \resizebox{\linewidth}{!}{\input{fxa.pgf}}
  \end{adjustwidth}
  \caption{f(x) aproximadas, más la original}
  \label{fig:aprox}
\end{figure}

\end{document}